% ─────────────────────────────────────────────────────────────────────────────
% SHARED coursework preamble — every course inputs THIS.
%
% PREAMBLE ONLY. No \begin{document}, no course-specific notation.
%
% Course-specific macros live at COURSE level, in
%   content/<course>/reference_docs/<course>_macros.tex
% which is inputted AFTER this file. ONE per course, shared by hw/qz/exam --
% macros under a single kind are unreachable from a sibling kind.
%
% That split is the point: the page style, the answer box, and the problem
% header are identical everywhere, so they are defined once; \conv and \ustep
% mean something in DSP and nothing in a networks course, so they are not.
%
% A per-problem file therefore starts:
%   % ─────────────────────────────────────────────────────────────────────────────
% SHARED coursework preamble — every course inputs THIS.
%
% PREAMBLE ONLY. No \begin{document}, no course-specific notation.
%
% Course-specific macros live at COURSE level, in
%   content/<course>/reference_docs/<course>_macros.tex
% which is inputted AFTER this file. ONE per course, shared by hw/qz/exam --
% macros under a single kind are unreachable from a sibling kind.
%
% That split is the point: the page style, the answer box, and the problem
% header are identical everywhere, so they are defined once; \conv and \ustep
% mean something in DSP and nothing in a networks course, so they are not.
%
% A per-problem file therefore starts:
%   % ─────────────────────────────────────────────────────────────────────────────
% SHARED coursework preamble — every course inputs THIS.
%
% PREAMBLE ONLY. No \begin{document}, no course-specific notation.
%
% Course-specific macros live at COURSE level, in
%   content/<course>/reference_docs/<course>_macros.tex
% which is inputted AFTER this file. ONE per course, shared by hw/qz/exam --
% macros under a single kind are unreachable from a sibling kind.
%
% That split is the point: the page style, the answer box, and the problem
% header are identical everywhere, so they are defined once; \conv and \ustep
% mean something in DSP and nothing in a networks course, so they are not.
%
% A per-problem file therefore starts:
%   \input{../../../../docs/latex/coursework_preamble.tex}
%   \input{../../reference_docs/cesc470_macros.tex}
%   \begin{document}
%
% Build-check one file, or a whole assignment:
%   docs/latex/build_tex.sh content/cesc_470/hw/hw01/p01_five_components.tex
%   docs/latex/build_tex.sh content/cesc_470/hw/hw01
%
% Doctrine: docs/directives/coursework-solutions.md
% ─────────────────────────────────────────────────────────────────────────────

\documentclass[11pt]{article}

\usepackage[margin=1in]{geometry}
\usepackage{amsmath, amssymb, mathtools}
\usepackage{graphicx}
\usepackage{float}
\usepackage{booktabs}
\usepackage{longtable}
\usepackage{enumitem}
\usepackage{siunitx}
\usepackage[dvipsnames]{xcolor}
\usepackage{tcolorbox}
\usepackage{fancyhdr}
\usepackage{caption}
\usepackage{tikz}
\usepackage{pgfplots}
\pgfplotsset{compat=1.18}
\usepackage[colorlinks=true, allcolors=NavyBlue]{hyperref}

\captionsetup{font=small, labelfont=bf}
\setlist[enumerate]{itemsep=2pt, topsep=4pt}
\setlist[itemize]{itemsep=2pt, topsep=4pt}

% --- final-answer box -------------------------------------------------------
% Every problem file ends with its result in one of these. The condensed
% solutions document is assembled by lifting exactly these.
\newtcolorbox{answerbox}[1][]{
  colback=Green!6, colframe=Green!55!black,
  boxrule=0.6pt, arc=2pt, left=6pt, right=6pt, top=4pt, bottom=4pt,
  #1
}

% --- citation to course material --------------------------------------------
% \cite{Module 01, PDF p55}  ->  a small italic parenthetical.
% Solutions cite the SLIDE/PAGE a formula came from so a grader (or a future
% reader) can check the claim against what was actually taught. See the
% directive: every non-obvious step carries one.
\newcommand{\srcref}[1]{{\small\itshape(#1)}}

% --- a step that came from OUTSIDE the course materials ---------------------
% Used only where the course is genuinely silent and the problem asks for
% outside research. Marked visually so it is never mistaken for lecture content.
\newcommand{\extref}[1]{{\small\itshape\color{Sepia}[external: #1]}}

% --- callout for the trap in a problem --------------------------------------
\newtcolorbox{trap}[1][]{
  colback=BurntOrange!7, colframe=BurntOrange!70!black,
  boxrule=0.6pt, arc=2pt, left=6pt, right=6pt, top=4pt, bottom=4pt,
  fonttitle=\bfseries, title=Watch out, #1
}

% --- per-problem header -----------------------------------------------------
% \problemheader{8}{15 pts}{Target clock rate}
% The optional 4-argument form carries a learning outcome where the course
% states one (CESC 410 does; CESC 470's HW1 does not).
\newcommand{\problemheader}[3]{%
  \begin{center}\large\bfseries
    Problem #1 \quad\normalsize\mdseries(#2)\\[2pt]
    \large\bfseries #3
  \end{center}
  \vspace{2pt}\hrule\vspace{8pt}
}
\newcommand{\problemheaderlo}[4]{%
  \begin{center}\large\bfseries
    Problem #1 \quad\normalsize\mdseries(#2, #3)\\[2pt]
    \large\bfseries #4
  \end{center}
  \vspace{2pt}\hrule\vspace{8pt}
}

% Course name in the running head. Defined BEFORE \lhead references it, and
% \renewcommand'd by each course's macros file (inputted after this one).
\providecommand{\coursename}{}

\pagestyle{fancy}
\fancyhf{}
\lhead{\small\coursename}
\rhead{\small\thepage}
\renewcommand{\headrulewidth}{0.3pt}

%   % CESC 470 Computer Architecture — course-specific macros ONLY.
%
% Inputted AFTER docs/latex/coursework_preamble.tex, which carries everything
% shared (answerbox, problemheader, srcref, page style). Put nothing general
% here: if another course would want it, it belongs in the shared preamble.

\renewcommand{\coursename}{CESC 470}

% --- performance-equation notation ------------------------------------------
% The course writes these in words on the slides; these render them compactly
% and, more importantly, IDENTICALLY across every problem file.
\newcommand{\CPUtime}{T_{\mathrm{CPU}}}
\newcommand{\IC}{\mathrm{IC}}                    % instruction count
\newcommand{\CPI}{\mathrm{CPI}}
\newcommand{\cycletime}{t_{\mathrm{cycle}}}
\newcommand{\clockrate}{f_{\mathrm{clk}}}
\newcommand{\cycles}{N_{\mathrm{cycles}}}
\newcommand{\speedup}{S}

% Amdahl's law, in the form the slides state it (Module 01, PDF p86, slide 49).
\newcommand{\amdahl}[3]{#1 + \dfrac{#2}{#3}}

%   \begin{document}
%
% Build-check one file, or a whole assignment:
%   docs/latex/build_tex.sh content/cesc_470/hw/hw01/p01_five_components.tex
%   docs/latex/build_tex.sh content/cesc_470/hw/hw01
%
% Doctrine: docs/directives/coursework-solutions.md
% ─────────────────────────────────────────────────────────────────────────────

\documentclass[11pt]{article}

\usepackage[margin=1in]{geometry}
\usepackage{amsmath, amssymb, mathtools}
\usepackage{graphicx}
\usepackage{float}
\usepackage{booktabs}
\usepackage{longtable}
\usepackage{enumitem}
\usepackage{siunitx}
\usepackage[dvipsnames]{xcolor}
\usepackage{tcolorbox}
\usepackage{fancyhdr}
\usepackage{caption}
\usepackage{tikz}
\usepackage{pgfplots}
\pgfplotsset{compat=1.18}
\usepackage[colorlinks=true, allcolors=NavyBlue]{hyperref}

\captionsetup{font=small, labelfont=bf}
\setlist[enumerate]{itemsep=2pt, topsep=4pt}
\setlist[itemize]{itemsep=2pt, topsep=4pt}

% --- final-answer box -------------------------------------------------------
% Every problem file ends with its result in one of these. The condensed
% solutions document is assembled by lifting exactly these.
\newtcolorbox{answerbox}[1][]{
  colback=Green!6, colframe=Green!55!black,
  boxrule=0.6pt, arc=2pt, left=6pt, right=6pt, top=4pt, bottom=4pt,
  #1
}

% --- citation to course material --------------------------------------------
% \cite{Module 01, PDF p55}  ->  a small italic parenthetical.
% Solutions cite the SLIDE/PAGE a formula came from so a grader (or a future
% reader) can check the claim against what was actually taught. See the
% directive: every non-obvious step carries one.
\newcommand{\srcref}[1]{{\small\itshape(#1)}}

% --- a step that came from OUTSIDE the course materials ---------------------
% Used only where the course is genuinely silent and the problem asks for
% outside research. Marked visually so it is never mistaken for lecture content.
\newcommand{\extref}[1]{{\small\itshape\color{Sepia}[external: #1]}}

% --- callout for the trap in a problem --------------------------------------
\newtcolorbox{trap}[1][]{
  colback=BurntOrange!7, colframe=BurntOrange!70!black,
  boxrule=0.6pt, arc=2pt, left=6pt, right=6pt, top=4pt, bottom=4pt,
  fonttitle=\bfseries, title=Watch out, #1
}

% --- per-problem header -----------------------------------------------------
% \problemheader{8}{15 pts}{Target clock rate}
% The optional 4-argument form carries a learning outcome where the course
% states one (CESC 410 does; CESC 470's HW1 does not).
\newcommand{\problemheader}[3]{%
  \begin{center}\large\bfseries
    Problem #1 \quad\normalsize\mdseries(#2)\\[2pt]
    \large\bfseries #3
  \end{center}
  \vspace{2pt}\hrule\vspace{8pt}
}
\newcommand{\problemheaderlo}[4]{%
  \begin{center}\large\bfseries
    Problem #1 \quad\normalsize\mdseries(#2, #3)\\[2pt]
    \large\bfseries #4
  \end{center}
  \vspace{2pt}\hrule\vspace{8pt}
}

% Course name in the running head. Defined BEFORE \lhead references it, and
% \renewcommand'd by each course's macros file (inputted after this one).
\providecommand{\coursename}{}

\pagestyle{fancy}
\fancyhf{}
\lhead{\small\coursename}
\rhead{\small\thepage}
\renewcommand{\headrulewidth}{0.3pt}

%   % CESC 470 Computer Architecture — course-specific macros ONLY.
%
% Inputted AFTER docs/latex/coursework_preamble.tex, which carries everything
% shared (answerbox, problemheader, srcref, page style). Put nothing general
% here: if another course would want it, it belongs in the shared preamble.

\renewcommand{\coursename}{CESC 470}

% --- performance-equation notation ------------------------------------------
% The course writes these in words on the slides; these render them compactly
% and, more importantly, IDENTICALLY across every problem file.
\newcommand{\CPUtime}{T_{\mathrm{CPU}}}
\newcommand{\IC}{\mathrm{IC}}                    % instruction count
\newcommand{\CPI}{\mathrm{CPI}}
\newcommand{\cycletime}{t_{\mathrm{cycle}}}
\newcommand{\clockrate}{f_{\mathrm{clk}}}
\newcommand{\cycles}{N_{\mathrm{cycles}}}
\newcommand{\speedup}{S}

% Amdahl's law, in the form the slides state it (Module 01, PDF p86, slide 49).
\newcommand{\amdahl}[3]{#1 + \dfrac{#2}{#3}}

%   \begin{document}
%
% Build-check one file, or a whole assignment:
%   docs/latex/build_tex.sh content/cesc_470/hw/hw01/p01_five_components.tex
%   docs/latex/build_tex.sh content/cesc_470/hw/hw01
%
% Doctrine: docs/directives/coursework-solutions.md
% ─────────────────────────────────────────────────────────────────────────────

\documentclass[11pt]{article}

\usepackage[margin=1in]{geometry}
\usepackage{amsmath, amssymb, mathtools}
\usepackage{graphicx}
\usepackage{float}
\usepackage{booktabs}
\usepackage{longtable}
\usepackage{enumitem}
\usepackage{siunitx}
\usepackage[dvipsnames]{xcolor}
\usepackage{tcolorbox}
\usepackage{fancyhdr}
\usepackage{caption}
\usepackage{tikz}
\usepackage{pgfplots}
\pgfplotsset{compat=1.18}
\usepackage[colorlinks=true, allcolors=NavyBlue]{hyperref}

\captionsetup{font=small, labelfont=bf}
\setlist[enumerate]{itemsep=2pt, topsep=4pt}
\setlist[itemize]{itemsep=2pt, topsep=4pt}

% --- final-answer box -------------------------------------------------------
% Every problem file ends with its result in one of these. The condensed
% solutions document is assembled by lifting exactly these.
\newtcolorbox{answerbox}[1][]{
  colback=Green!6, colframe=Green!55!black,
  boxrule=0.6pt, arc=2pt, left=6pt, right=6pt, top=4pt, bottom=4pt,
  #1
}

% --- citation to course material --------------------------------------------
% \cite{Module 01, PDF p55}  ->  a small italic parenthetical.
% Solutions cite the SLIDE/PAGE a formula came from so a grader (or a future
% reader) can check the claim against what was actually taught. See the
% directive: every non-obvious step carries one.
\newcommand{\srcref}[1]{{\small\itshape(#1)}}

% --- a step that came from OUTSIDE the course materials ---------------------
% Used only where the course is genuinely silent and the problem asks for
% outside research. Marked visually so it is never mistaken for lecture content.
\newcommand{\extref}[1]{{\small\itshape\color{Sepia}[external: #1]}}

% --- callout for the trap in a problem --------------------------------------
\newtcolorbox{trap}[1][]{
  colback=BurntOrange!7, colframe=BurntOrange!70!black,
  boxrule=0.6pt, arc=2pt, left=6pt, right=6pt, top=4pt, bottom=4pt,
  fonttitle=\bfseries, title=Watch out, #1
}

% --- per-problem header -----------------------------------------------------
% \problemheader{8}{15 pts}{Target clock rate}
% The optional 4-argument form carries a learning outcome where the course
% states one (CESC 410 does; CESC 470's HW1 does not).
\newcommand{\problemheader}[3]{%
  \begin{center}\large\bfseries
    Problem #1 \quad\normalsize\mdseries(#2)\\[2pt]
    \large\bfseries #3
  \end{center}
  \vspace{2pt}\hrule\vspace{8pt}
}
\newcommand{\problemheaderlo}[4]{%
  \begin{center}\large\bfseries
    Problem #1 \quad\normalsize\mdseries(#2, #3)\\[2pt]
    \large\bfseries #4
  \end{center}
  \vspace{2pt}\hrule\vspace{8pt}
}

% Course name in the running head. Defined BEFORE \lhead references it, and
% \renewcommand'd by each course's macros file (inputted after this one).
\providecommand{\coursename}{}

\pagestyle{fancy}
\fancyhf{}
\lhead{\small\coursename}
\rhead{\small\thepage}
\renewcommand{\headrulewidth}{0.3pt}

% ─────────────────────────────────────────────────────────────────────────────
% CESC 410 course-specific macros. Inputted AFTER the shared preamble:
%
%   % ─────────────────────────────────────────────────────────────────────────────
% SHARED coursework preamble — every course inputs THIS.
%
% PREAMBLE ONLY. No \begin{document}, no course-specific notation.
%
% Course-specific macros live at COURSE level, in
%   content/<course>/reference_docs/<course>_macros.tex
% which is inputted AFTER this file. ONE per course, shared by hw/qz/exam --
% macros under a single kind are unreachable from a sibling kind.
%
% That split is the point: the page style, the answer box, and the problem
% header are identical everywhere, so they are defined once; \conv and \ustep
% mean something in DSP and nothing in a networks course, so they are not.
%
% A per-problem file therefore starts:
%   % ─────────────────────────────────────────────────────────────────────────────
% SHARED coursework preamble — every course inputs THIS.
%
% PREAMBLE ONLY. No \begin{document}, no course-specific notation.
%
% Course-specific macros live at COURSE level, in
%   content/<course>/reference_docs/<course>_macros.tex
% which is inputted AFTER this file. ONE per course, shared by hw/qz/exam --
% macros under a single kind are unreachable from a sibling kind.
%
% That split is the point: the page style, the answer box, and the problem
% header are identical everywhere, so they are defined once; \conv and \ustep
% mean something in DSP and nothing in a networks course, so they are not.
%
% A per-problem file therefore starts:
%   \input{../../../../docs/latex/coursework_preamble.tex}
%   \input{../../reference_docs/cesc470_macros.tex}
%   \begin{document}
%
% Build-check one file, or a whole assignment:
%   docs/latex/build_tex.sh content/cesc_470/hw/hw01/p01_five_components.tex
%   docs/latex/build_tex.sh content/cesc_470/hw/hw01
%
% Doctrine: docs/directives/coursework-solutions.md
% ─────────────────────────────────────────────────────────────────────────────

\documentclass[11pt]{article}

\usepackage[margin=1in]{geometry}
\usepackage{amsmath, amssymb, mathtools}
\usepackage{graphicx}
\usepackage{float}
\usepackage{booktabs}
\usepackage{longtable}
\usepackage{enumitem}
\usepackage{siunitx}
\usepackage[dvipsnames]{xcolor}
\usepackage{tcolorbox}
\usepackage{fancyhdr}
\usepackage{caption}
\usepackage{tikz}
\usepackage{pgfplots}
\pgfplotsset{compat=1.18}
\usepackage[colorlinks=true, allcolors=NavyBlue]{hyperref}

\captionsetup{font=small, labelfont=bf}
\setlist[enumerate]{itemsep=2pt, topsep=4pt}
\setlist[itemize]{itemsep=2pt, topsep=4pt}

% --- final-answer box -------------------------------------------------------
% Every problem file ends with its result in one of these. The condensed
% solutions document is assembled by lifting exactly these.
\newtcolorbox{answerbox}[1][]{
  colback=Green!6, colframe=Green!55!black,
  boxrule=0.6pt, arc=2pt, left=6pt, right=6pt, top=4pt, bottom=4pt,
  #1
}

% --- citation to course material --------------------------------------------
% \cite{Module 01, PDF p55}  ->  a small italic parenthetical.
% Solutions cite the SLIDE/PAGE a formula came from so a grader (or a future
% reader) can check the claim against what was actually taught. See the
% directive: every non-obvious step carries one.
\newcommand{\srcref}[1]{{\small\itshape(#1)}}

% --- a step that came from OUTSIDE the course materials ---------------------
% Used only where the course is genuinely silent and the problem asks for
% outside research. Marked visually so it is never mistaken for lecture content.
\newcommand{\extref}[1]{{\small\itshape\color{Sepia}[external: #1]}}

% --- callout for the trap in a problem --------------------------------------
\newtcolorbox{trap}[1][]{
  colback=BurntOrange!7, colframe=BurntOrange!70!black,
  boxrule=0.6pt, arc=2pt, left=6pt, right=6pt, top=4pt, bottom=4pt,
  fonttitle=\bfseries, title=Watch out, #1
}

% --- per-problem header -----------------------------------------------------
% \problemheader{8}{15 pts}{Target clock rate}
% The optional 4-argument form carries a learning outcome where the course
% states one (CESC 410 does; CESC 470's HW1 does not).
\newcommand{\problemheader}[3]{%
  \begin{center}\large\bfseries
    Problem #1 \quad\normalsize\mdseries(#2)\\[2pt]
    \large\bfseries #3
  \end{center}
  \vspace{2pt}\hrule\vspace{8pt}
}
\newcommand{\problemheaderlo}[4]{%
  \begin{center}\large\bfseries
    Problem #1 \quad\normalsize\mdseries(#2, #3)\\[2pt]
    \large\bfseries #4
  \end{center}
  \vspace{2pt}\hrule\vspace{8pt}
}

% Course name in the running head. Defined BEFORE \lhead references it, and
% \renewcommand'd by each course's macros file (inputted after this one).
\providecommand{\coursename}{}

\pagestyle{fancy}
\fancyhf{}
\lhead{\small\coursename}
\rhead{\small\thepage}
\renewcommand{\headrulewidth}{0.3pt}

%   % CESC 470 Computer Architecture — course-specific macros ONLY.
%
% Inputted AFTER docs/latex/coursework_preamble.tex, which carries everything
% shared (answerbox, problemheader, srcref, page style). Put nothing general
% here: if another course would want it, it belongs in the shared preamble.

\renewcommand{\coursename}{CESC 470}

% --- performance-equation notation ------------------------------------------
% The course writes these in words on the slides; these render them compactly
% and, more importantly, IDENTICALLY across every problem file.
\newcommand{\CPUtime}{T_{\mathrm{CPU}}}
\newcommand{\IC}{\mathrm{IC}}                    % instruction count
\newcommand{\CPI}{\mathrm{CPI}}
\newcommand{\cycletime}{t_{\mathrm{cycle}}}
\newcommand{\clockrate}{f_{\mathrm{clk}}}
\newcommand{\cycles}{N_{\mathrm{cycles}}}
\newcommand{\speedup}{S}

% Amdahl's law, in the form the slides state it (Module 01, PDF p86, slide 49).
\newcommand{\amdahl}[3]{#1 + \dfrac{#2}{#3}}

%   \begin{document}
%
% Build-check one file, or a whole assignment:
%   docs/latex/build_tex.sh content/cesc_470/hw/hw01/p01_five_components.tex
%   docs/latex/build_tex.sh content/cesc_470/hw/hw01
%
% Doctrine: docs/directives/coursework-solutions.md
% ─────────────────────────────────────────────────────────────────────────────

\documentclass[11pt]{article}

\usepackage[margin=1in]{geometry}
\usepackage{amsmath, amssymb, mathtools}
\usepackage{graphicx}
\usepackage{float}
\usepackage{booktabs}
\usepackage{longtable}
\usepackage{enumitem}
\usepackage{siunitx}
\usepackage[dvipsnames]{xcolor}
\usepackage{tcolorbox}
\usepackage{fancyhdr}
\usepackage{caption}
\usepackage{tikz}
\usepackage{pgfplots}
\pgfplotsset{compat=1.18}
\usepackage[colorlinks=true, allcolors=NavyBlue]{hyperref}

\captionsetup{font=small, labelfont=bf}
\setlist[enumerate]{itemsep=2pt, topsep=4pt}
\setlist[itemize]{itemsep=2pt, topsep=4pt}

% --- final-answer box -------------------------------------------------------
% Every problem file ends with its result in one of these. The condensed
% solutions document is assembled by lifting exactly these.
\newtcolorbox{answerbox}[1][]{
  colback=Green!6, colframe=Green!55!black,
  boxrule=0.6pt, arc=2pt, left=6pt, right=6pt, top=4pt, bottom=4pt,
  #1
}

% --- citation to course material --------------------------------------------
% \cite{Module 01, PDF p55}  ->  a small italic parenthetical.
% Solutions cite the SLIDE/PAGE a formula came from so a grader (or a future
% reader) can check the claim against what was actually taught. See the
% directive: every non-obvious step carries one.
\newcommand{\srcref}[1]{{\small\itshape(#1)}}

% --- a step that came from OUTSIDE the course materials ---------------------
% Used only where the course is genuinely silent and the problem asks for
% outside research. Marked visually so it is never mistaken for lecture content.
\newcommand{\extref}[1]{{\small\itshape\color{Sepia}[external: #1]}}

% --- callout for the trap in a problem --------------------------------------
\newtcolorbox{trap}[1][]{
  colback=BurntOrange!7, colframe=BurntOrange!70!black,
  boxrule=0.6pt, arc=2pt, left=6pt, right=6pt, top=4pt, bottom=4pt,
  fonttitle=\bfseries, title=Watch out, #1
}

% --- per-problem header -----------------------------------------------------
% \problemheader{8}{15 pts}{Target clock rate}
% The optional 4-argument form carries a learning outcome where the course
% states one (CESC 410 does; CESC 470's HW1 does not).
\newcommand{\problemheader}[3]{%
  \begin{center}\large\bfseries
    Problem #1 \quad\normalsize\mdseries(#2)\\[2pt]
    \large\bfseries #3
  \end{center}
  \vspace{2pt}\hrule\vspace{8pt}
}
\newcommand{\problemheaderlo}[4]{%
  \begin{center}\large\bfseries
    Problem #1 \quad\normalsize\mdseries(#2, #3)\\[2pt]
    \large\bfseries #4
  \end{center}
  \vspace{2pt}\hrule\vspace{8pt}
}

% Course name in the running head. Defined BEFORE \lhead references it, and
% \renewcommand'd by each course's macros file (inputted after this one).
\providecommand{\coursename}{}

\pagestyle{fancy}
\fancyhf{}
\lhead{\small\coursename}
\rhead{\small\thepage}
\renewcommand{\headrulewidth}{0.3pt}

%   % ─────────────────────────────────────────────────────────────────────────────
% CESC 410 course-specific macros. Inputted AFTER the shared preamble:
%
%   % ─────────────────────────────────────────────────────────────────────────────
% SHARED coursework preamble — every course inputs THIS.
%
% PREAMBLE ONLY. No \begin{document}, no course-specific notation.
%
% Course-specific macros live at COURSE level, in
%   content/<course>/reference_docs/<course>_macros.tex
% which is inputted AFTER this file. ONE per course, shared by hw/qz/exam --
% macros under a single kind are unreachable from a sibling kind.
%
% That split is the point: the page style, the answer box, and the problem
% header are identical everywhere, so they are defined once; \conv and \ustep
% mean something in DSP and nothing in a networks course, so they are not.
%
% A per-problem file therefore starts:
%   \input{../../../../docs/latex/coursework_preamble.tex}
%   \input{../../reference_docs/cesc470_macros.tex}
%   \begin{document}
%
% Build-check one file, or a whole assignment:
%   docs/latex/build_tex.sh content/cesc_470/hw/hw01/p01_five_components.tex
%   docs/latex/build_tex.sh content/cesc_470/hw/hw01
%
% Doctrine: docs/directives/coursework-solutions.md
% ─────────────────────────────────────────────────────────────────────────────

\documentclass[11pt]{article}

\usepackage[margin=1in]{geometry}
\usepackage{amsmath, amssymb, mathtools}
\usepackage{graphicx}
\usepackage{float}
\usepackage{booktabs}
\usepackage{longtable}
\usepackage{enumitem}
\usepackage{siunitx}
\usepackage[dvipsnames]{xcolor}
\usepackage{tcolorbox}
\usepackage{fancyhdr}
\usepackage{caption}
\usepackage{tikz}
\usepackage{pgfplots}
\pgfplotsset{compat=1.18}
\usepackage[colorlinks=true, allcolors=NavyBlue]{hyperref}

\captionsetup{font=small, labelfont=bf}
\setlist[enumerate]{itemsep=2pt, topsep=4pt}
\setlist[itemize]{itemsep=2pt, topsep=4pt}

% --- final-answer box -------------------------------------------------------
% Every problem file ends with its result in one of these. The condensed
% solutions document is assembled by lifting exactly these.
\newtcolorbox{answerbox}[1][]{
  colback=Green!6, colframe=Green!55!black,
  boxrule=0.6pt, arc=2pt, left=6pt, right=6pt, top=4pt, bottom=4pt,
  #1
}

% --- citation to course material --------------------------------------------
% \cite{Module 01, PDF p55}  ->  a small italic parenthetical.
% Solutions cite the SLIDE/PAGE a formula came from so a grader (or a future
% reader) can check the claim against what was actually taught. See the
% directive: every non-obvious step carries one.
\newcommand{\srcref}[1]{{\small\itshape(#1)}}

% --- a step that came from OUTSIDE the course materials ---------------------
% Used only where the course is genuinely silent and the problem asks for
% outside research. Marked visually so it is never mistaken for lecture content.
\newcommand{\extref}[1]{{\small\itshape\color{Sepia}[external: #1]}}

% --- callout for the trap in a problem --------------------------------------
\newtcolorbox{trap}[1][]{
  colback=BurntOrange!7, colframe=BurntOrange!70!black,
  boxrule=0.6pt, arc=2pt, left=6pt, right=6pt, top=4pt, bottom=4pt,
  fonttitle=\bfseries, title=Watch out, #1
}

% --- per-problem header -----------------------------------------------------
% \problemheader{8}{15 pts}{Target clock rate}
% The optional 4-argument form carries a learning outcome where the course
% states one (CESC 410 does; CESC 470's HW1 does not).
\newcommand{\problemheader}[3]{%
  \begin{center}\large\bfseries
    Problem #1 \quad\normalsize\mdseries(#2)\\[2pt]
    \large\bfseries #3
  \end{center}
  \vspace{2pt}\hrule\vspace{8pt}
}
\newcommand{\problemheaderlo}[4]{%
  \begin{center}\large\bfseries
    Problem #1 \quad\normalsize\mdseries(#2, #3)\\[2pt]
    \large\bfseries #4
  \end{center}
  \vspace{2pt}\hrule\vspace{8pt}
}

% Course name in the running head. Defined BEFORE \lhead references it, and
% \renewcommand'd by each course's macros file (inputted after this one).
\providecommand{\coursename}{}

\pagestyle{fancy}
\fancyhf{}
\lhead{\small\coursename}
\rhead{\small\thepage}
\renewcommand{\headrulewidth}{0.3pt}

%   % ─────────────────────────────────────────────────────────────────────────────
% CESC 410 course-specific macros. Inputted AFTER the shared preamble:
%
%   \input{../../../../docs/latex/coursework_preamble.tex}
%   \input{../reference_docs/cesc410_macros.tex}
%   \begin{document}
%
% ONLY things that mean something in DSP go here. The document class, the
% package set, \answerbox, \srcref, \extref, the trap callout, \problemheader
% and the page style are identical across courses and live in the shared
% preamble — do not redefine them here.
%
% Build-check one problem in isolation, or the whole assignment:
%   docs/latex/build_tex.sh content/cesc_410/hw/hw01/p01_phasor_form.tex
%   docs/latex/build_tex.sh content/cesc_410/hw/hw01
% ─────────────────────────────────────────────────────────────────────────────

% Fills the running head defined by the shared preamble.
\renewcommand{\coursename}{CESC 410}

% --- DSP notation -----------------------------------------------------------
% Defined once so every problem writes the same symbols. The course uses a
% cursive w for omega on the board; we render omega throughout.
% DTFT. Takes the SYMBOL so \dtft{Y} gives Y(e^{jw}); previously it declared an
% argument and then ignored it, so \dtft{Y} silently printed X(e^{jw}) -- the
% wrong signal, with no error. Unused in HW1, so nothing rendered was affected.
\newcommand{\dtft}[1]{#1\!\left(e^{j\omega}\right)}
\newcommand{\ztr}[1]{\mathcal{Z}\!\left\{#1\right\}}    % z-transform
\newcommand{\conv}{\ast}                                % convolution
\newcommand{\Real}{\operatorname{Re}}
\newcommand{\Imag}{\operatorname{Im}}
% Magnitude / angle. The \mathord wrapper matters: \angle is a mathord, so a
% following "-" would be parsed as a BINARY minus and typeset as "10 /_ - 75",
% with spacing on both sides. Wrapping starts a fresh math list, making the
% sign unary: "10 /_ -75".
\newcommand{\phasor}[2]{#1\,\angle\,\mathord{#2}}
\newcommand{\ustep}{u}                                  % unit step u[n]
\DeclareMathOperator{\sinc}{sinc}

% --- stem plot helper -------------------------------------------------------
% \stemplot{-2}{5}{{0,1},{1,-0.7},{2,0.5}}   % xmin, xmax, {n,value} pairs
% Matches the plotting format the lecture notes use: open circles on stems,
% zero line drawn through.
\newcommand{\stemplot}[4][]{%
  \begin{tikzpicture}[baseline]
    \begin{axis}[
      width=0.62\textwidth, height=4.4cm,
      xmin=#2-0.6, xmax=#3+0.6,
      axis lines=middle, xlabel={$n$}, ylabel={$x[n]$},
      % Park the y-label at the top-left corner of the axis box. With
      % `axis lines=middle` its default position is mid-height ON the y-axis,
      % where a full-height stem at n=0 runs straight through it.
      ylabel style={at={(current axis.north west)}, anchor=south west},
      xtick={#2,...,#3}, tick label style={font=\scriptsize},
      ymajorgrids=false, enlarge y limits=0.25,
      #1
    ]
      \addplot+[ycomb, mark=o, mark options={fill=white}, thick] coordinates {#4};
    \end{axis}
  \end{tikzpicture}%
}

%   \begin{document}
%
% ONLY things that mean something in DSP go here. The document class, the
% package set, \answerbox, \srcref, \extref, the trap callout, \problemheader
% and the page style are identical across courses and live in the shared
% preamble — do not redefine them here.
%
% Build-check one problem in isolation, or the whole assignment:
%   docs/latex/build_tex.sh content/cesc_410/hw/hw01/p01_phasor_form.tex
%   docs/latex/build_tex.sh content/cesc_410/hw/hw01
% ─────────────────────────────────────────────────────────────────────────────

% Fills the running head defined by the shared preamble.
\renewcommand{\coursename}{CESC 410}

% --- DSP notation -----------------------------------------------------------
% Defined once so every problem writes the same symbols. The course uses a
% cursive w for omega on the board; we render omega throughout.
% DTFT. Takes the SYMBOL so \dtft{Y} gives Y(e^{jw}); previously it declared an
% argument and then ignored it, so \dtft{Y} silently printed X(e^{jw}) -- the
% wrong signal, with no error. Unused in HW1, so nothing rendered was affected.
\newcommand{\dtft}[1]{#1\!\left(e^{j\omega}\right)}
\newcommand{\ztr}[1]{\mathcal{Z}\!\left\{#1\right\}}    % z-transform
\newcommand{\conv}{\ast}                                % convolution
\newcommand{\Real}{\operatorname{Re}}
\newcommand{\Imag}{\operatorname{Im}}
% Magnitude / angle. The \mathord wrapper matters: \angle is a mathord, so a
% following "-" would be parsed as a BINARY minus and typeset as "10 /_ - 75",
% with spacing on both sides. Wrapping starts a fresh math list, making the
% sign unary: "10 /_ -75".
\newcommand{\phasor}[2]{#1\,\angle\,\mathord{#2}}
\newcommand{\ustep}{u}                                  % unit step u[n]
\DeclareMathOperator{\sinc}{sinc}

% --- stem plot helper -------------------------------------------------------
% \stemplot{-2}{5}{{0,1},{1,-0.7},{2,0.5}}   % xmin, xmax, {n,value} pairs
% Matches the plotting format the lecture notes use: open circles on stems,
% zero line drawn through.
\newcommand{\stemplot}[4][]{%
  \begin{tikzpicture}[baseline]
    \begin{axis}[
      width=0.62\textwidth, height=4.4cm,
      xmin=#2-0.6, xmax=#3+0.6,
      axis lines=middle, xlabel={$n$}, ylabel={$x[n]$},
      % Park the y-label at the top-left corner of the axis box. With
      % `axis lines=middle` its default position is mid-height ON the y-axis,
      % where a full-height stem at n=0 runs straight through it.
      ylabel style={at={(current axis.north west)}, anchor=south west},
      xtick={#2,...,#3}, tick label style={font=\scriptsize},
      ymajorgrids=false, enlarge y limits=0.25,
      #1
    ]
      \addplot+[ycomb, mark=o, mark options={fill=white}, thick] coordinates {#4};
    \end{axis}
  \end{tikzpicture}%
}

%   \begin{document}
%
% ONLY things that mean something in DSP go here. The document class, the
% package set, \answerbox, \srcref, \extref, the trap callout, \problemheader
% and the page style are identical across courses and live in the shared
% preamble — do not redefine them here.
%
% Build-check one problem in isolation, or the whole assignment:
%   docs/latex/build_tex.sh content/cesc_410/hw/hw01/p01_phasor_form.tex
%   docs/latex/build_tex.sh content/cesc_410/hw/hw01
% ─────────────────────────────────────────────────────────────────────────────

% Fills the running head defined by the shared preamble.
\renewcommand{\coursename}{CESC 410}

% --- DSP notation -----------------------------------------------------------
% Defined once so every problem writes the same symbols. The course uses a
% cursive w for omega on the board; we render omega throughout.
% DTFT. Takes the SYMBOL so \dtft{Y} gives Y(e^{jw}); previously it declared an
% argument and then ignored it, so \dtft{Y} silently printed X(e^{jw}) -- the
% wrong signal, with no error. Unused in HW1, so nothing rendered was affected.
\newcommand{\dtft}[1]{#1\!\left(e^{j\omega}\right)}
\newcommand{\ztr}[1]{\mathcal{Z}\!\left\{#1\right\}}    % z-transform
\newcommand{\conv}{\ast}                                % convolution
\newcommand{\Real}{\operatorname{Re}}
\newcommand{\Imag}{\operatorname{Im}}
% Magnitude / angle. The \mathord wrapper matters: \angle is a mathord, so a
% following "-" would be parsed as a BINARY minus and typeset as "10 /_ - 75",
% with spacing on both sides. Wrapping starts a fresh math list, making the
% sign unary: "10 /_ -75".
\newcommand{\phasor}[2]{#1\,\angle\,\mathord{#2}}
\newcommand{\ustep}{u}                                  % unit step u[n]
\DeclareMathOperator{\sinc}{sinc}

% --- stem plot helper -------------------------------------------------------
% \stemplot{-2}{5}{{0,1},{1,-0.7},{2,0.5}}   % xmin, xmax, {n,value} pairs
% Matches the plotting format the lecture notes use: open circles on stems,
% zero line drawn through.
\newcommand{\stemplot}[4][]{%
  \begin{tikzpicture}[baseline]
    \begin{axis}[
      width=0.62\textwidth, height=4.4cm,
      xmin=#2-0.6, xmax=#3+0.6,
      axis lines=middle, xlabel={$n$}, ylabel={$x[n]$},
      % Park the y-label at the top-left corner of the axis box. With
      % `axis lines=middle` its default position is mid-height ON the y-axis,
      % where a full-height stem at n=0 runs straight through it.
      ylabel style={at={(current axis.north west)}, anchor=south west},
      xtick={#2,...,#3}, tick label style={font=\scriptsize},
      ymajorgrids=false, enlarge y limits=0.25,
      #1
    ]
      \addplot+[ycomb, mark=o, mark options={fill=white}, thick] coordinates {#4};
    \end{axis}
  \end{tikzpicture}%
}


% CONDENSED ANSWER SHEET. Final answers only -- no derivations.
% The full step-by-step work for each problem lives in its own file:
%   p01_phasor_form.tex  p02_phasor_arithmetic.tex  p03_signal_transformations.tex
%   p04_periodicity.tex  p05_system_properties.tex  p06_convolution.tex

\title{\vspace{-1.5cm}CESC 410 --- Homework 1 Solutions\\
       {\large Basic Signals/Systems and Convolution}}
\author{}
\date{Fall 2026}

\begin{document}
\maketitle
\vspace{-1.2cm}

%---------------------------------------------------------------------------
\section*{Problem 1 \normalsize(LO01, 20 pts) --- Phasor form}

\begin{align*}
  x_a(t) = 10\sin(\omega t + \pi/12)  &\;\Rightarrow\; \mathbf{X}_a = \phasor{10}{-5\pi/12} = \phasor{10}{-75^\circ} \\
  x_b(t) = 15\cos(\omega t + \pi/12)  &\;\Rightarrow\; \mathbf{X}_b = \phasor{15}{\pi/12}   = \phasor{15}{15^\circ}  \\
  x_c(t) = -12\sin(\omega t - \pi/12) &\;\Rightarrow\; \mathbf{X}_c = \phasor{12}{5\pi/12}  = \phasor{12}{75^\circ}  \\
  x_d(t) = -8\cos(\omega t - \pi/12)  &\;\Rightarrow\; \mathbf{X}_d = \phasor{8}{11\pi/12}  = \phasor{8}{165^\circ}
\end{align*}

%---------------------------------------------------------------------------
\section*{Problem 2 \normalsize(LO01, 10 pts) --- Phasor arithmetic}

\begin{align*}
  x_a(t) + x_b(t) &= \phasor{5\sqrt{13}}{-18.69^\circ} = \phasor{18.03}{-18.69^\circ}
    \;\Rightarrow\; 18.03\cos(\omega t - 18.69^\circ) \\
  x_b(t) - x_c(t) &= \phasor{3\sqrt{21}}{-34.11^\circ} = \phasor{13.75}{-34.11^\circ}
    \;\Rightarrow\; 13.75\cos(\omega t - 34.11^\circ)
\end{align*}

%---------------------------------------------------------------------------
\section*{Problem 3 \normalsize(LO02, 20 pts) --- Signal transformations}

Samples of $x[n] = \cos(\pi n/6)$ on $0 \le n \le 3$: \;$\{1,\; \sqrt{3}/2 \approx 0.866,\; 0.5,\; 0\}$.

\begin{center}
\begin{tabular}{llll}
  \toprule
  & Signal & Support & Samples (left to right) \\
  \midrule
  1 & $x[n]$    & $0 \le n \le 3$   & $1,\ 0.866,\ 0.5,\ 0$ \\
  2 & $x[n-2]$  & $2 \le n \le 5$   & $1,\ 0.866,\ 0.5,\ 0$ \\
  3 & $x[-n]$   & $-3 \le n \le 0$  & $0,\ 0.5,\ 0.866,\ 1$ \\
  4 & $x[2-n]$  & $-1 \le n \le 2$  & $0,\ 0.5,\ 0.866,\ 1$ \\
  \bottomrule
\end{tabular}
\end{center}

\vspace{-2pt}
\begin{center}
\begin{minipage}{0.48\textwidth}\centering
  {\footnotesize 1.\ $x[n]$}\\
  \stemplot[width=\textwidth,height=3.4cm]{-4}{6}{(-4,0)(-3,0)(-2,0)(-1,0)(0,1)(1,0.866)(2,0.5)(3,0)(4,0)(5,0)(6,0)}
\end{minipage}\hfill
\begin{minipage}{0.48\textwidth}\centering
  {\footnotesize 2.\ $x[n-2]$}\\
  \stemplot[width=\textwidth,height=3.4cm]{-4}{6}{(-4,0)(-3,0)(-2,0)(-1,0)(0,0)(1,0)(2,1)(3,0.866)(4,0.5)(5,0)(6,0)}
\end{minipage}

\vspace{4pt}
\begin{minipage}{0.48\textwidth}\centering
  {\footnotesize 3.\ $x[-n]$}\\
  \stemplot[width=\textwidth,height=3.4cm]{-6}{4}{(-6,0)(-5,0)(-4,0)(-3,0)(-2,0.5)(-1,0.866)(0,1)(1,0)(2,0)(3,0)(4,0)}
\end{minipage}\hfill
\begin{minipage}{0.48\textwidth}\centering
  {\footnotesize 4.\ $x[2-n]$}\\
  \stemplot[width=\textwidth,height=3.4cm]{-5}{5}{(-5,0)(-4,0)(-3,0)(-2,0)(-1,0)(0,0.5)(1,0.866)(2,1)(3,0)(4,0)(5,0)}
\end{minipage}
\end{center}

%---------------------------------------------------------------------------
\section*{Problem 4 \normalsize(LO03, 20 pts) --- Periodicity}

\begin{center}
\begin{tabular}{clllc}
  \toprule
  & Signal & $\omega_0/2\pi$ & Periodic? & $N_0$ \\
  \midrule
  1 & $x_a[n] = e^{j(\pi n/8)}$        & $1/16$        & Yes & $16$ \\
  2 & $x_b[n] = e^{j(7\pi n/9)}$       & $7/18$        & Yes & $18$ \\
  3 & $x_c[n] = n\,e^{j(\pi n/8)}$     & ---           & No (envelope $n$ grows) & --- \\
  4 & $x_d[n] = e^{j(\pi n/\sqrt{8})}$ & $\sqrt{2}/8$ (irrational) & No & --- \\
  \bottomrule
\end{tabular}
\end{center}

%---------------------------------------------------------------------------
\section*{Problem 5 \normalsize(LO04, 40 pts) --- System properties}

\begin{center}
\begin{tabular}{clcccc}
  \toprule
  & System & Linear & Time-invariant & Causal & BIBO stable \\
  \midrule
  1 & $T_a(x[n]) = n\,x[n]$          & Yes & No  & Yes & No  \\
  2 & $T_b(x[n]) = a\,x[n] + c$      & No$^\dagger$ & Yes & Yes & Yes \\
  3 & $T_c(x[n]) = x[n+1] + x[n-1]$  & Yes & Yes & No  & Yes \\
  4 & $T_d(x[n]) = x[n]\,x[n+1]$     & No  & Yes & No  & Yes \\
  5 & $T_e(x[n]) = e^{x[n-1]}$       & No  & Yes & Yes & Yes \\
  \bottomrule
\end{tabular}
\end{center}

{\footnotesize $^\dagger$ Affine, not linear, for $c \ne 0$. Linear if $c = 0$.}

%---------------------------------------------------------------------------
\section*{Problem 6 \normalsize(LO05, 20 pts) --- Convolution}

{\footnotesize\itshape (Labelled ``Prob 1'' a second time in the handout.)}

For $x[n] = a^n\ustep[n]$:
\begin{equation*}
  y[n] = x[n] \conv x[n]
       = \sum_{k=0}^{n} a^{k} a^{\,n-k}
       = \sum_{k=0}^{n} a^{n}
       = (n+1)\,a^{n}\,\ustep[n].
\end{equation*}

\end{document}
